%TODO: forall exists
%TODO: ranges
\chapter{Type Analysis for Symbolic Execution}
\label{ch:typ}

The main focus of this thesis is the analysis of \rego policies based on symbolic documents, i.e., documents where every field is specified only by a type and not a concrete value.
As it was already mentioned, the approach of this work is to encode given policies into SMT models and reason on them.

In Chapter~\ref{ch:form}, we defined $\typof$ as a substitution mapping terms to values in a concrete evaluation context.
This chapter extends the definition of $\typof$ to operate on symbolic documents using a set of inference rules.
Before these rules are provided, we refine the type system for the symbolic execution.

% To effectively and correctly create a SMT model of policy specified in Rego language, it is important to infer datatypes of appearing values, as precisely as possible.
% In this chapter, we describe a set of rules followed by the type analysis, which creates a function $\typof$, mapping every value to its corresponding type.

\section{Type System Extension}

During a symbolic execution, each term $t \in \Terms$ does not map to a single value, but rather a (possibly infinite) set of possible values $\{v_1,\cdots,v_n\}$, where $v_i \in \Vals$.
Given that \rego is dynamically typed, these values do not have to be of the same type.
This implies that we need to extend the definition of the type system to map terms to a set of all possible types.

\subsection{The Set of Types}

For symbolic execution, i.e., for the rest of this paper, we redefine the set of \rego types $\Types$ with the definition given by the grammar:
\begin{align*}
  \ty \;::=\; & \{\Str\} \mid \{\Int\} \mid \{\Bool\} \mid \{\Null\} \mid \{\Undef\} \mid \{\Unknown\} \\
   \mid\;   & \{\Array{\ty}\} \\
   \mid\;   & \{\RegoSet{\ty}\} \\
   \mid\;   & \{\Obj{v_1{:}\ty_1,\ldots,v_n{:}\ty_n} \\
   \mid\;   & \ty_1 \cup \ty_2
\end{align*}
where:
\begin{itemize}
  \item $\ty_1 \cup \ty_2$ is the standard set union,
  \item $\Array{\ty}$ denotes the array type with elements of type $\ty$,
  \item $\RegoSet{\ty}$ is a \rego set type of elements of type $\ty$,
  \item $\Obj{v_1{:}\ty_1,\ldots,v_n{:}\ty_n}$ is an object type with fields $v_i \in \Vals$ typed $\ty_i$,
  \item $\{\Unknown\}$ is a special \emph{placeholder} type used when no information is available.
\end{itemize}

For the rest of this chapter, we simplify the notation of given types and omit the wrapping curly braces for singleton sets.
With this notation, e.g., $\Str$ is equal to $\{\Str\}$.


% \section{Properties of Types}
%
% This section introduces additional properties of Rego types, that it their \emph{depth} and \emph{precision order}.
% Additionally, this section presents \emph{union normalization}, i.e., a formal process to transform the union type into a standard (normalized) form.
%
% TODO: move to transformation TODO TODO
% \subsection{Type depth}
%
% The \emph{depth} $\depth$ of a types is defined as follows:
% \begin{itemize}
%   \item $\depth(\ty) = 0$, for $\ty$ being an atomic type,
%   \item $\depth(\ty_1 \cup \cdots \cup \ty_m) = \depth(\ty_i)$, for $1 \leq i \leq m$, where $\forall_{1 \leq j \leq m}:\depth(\ty_i) \geq \depth(\ty_j)$,
%   \item $\depth(\Obj{f_1{:}\ty_1,\ldots,f_n{:}\ty_n}) = 1+\depth(\ty_1 \cup \cdots \cup \ty_n)$,
%   \item $\depth(\Array{\ty}) = 1+\depth(\ty)$,
%   \item $\depth(\RegoSet{\ty}) = 1+\depth(\ty)$.
% \end{itemize}
%
% Informally said, atomic types have depth $0$, and composite types have their depth greater by one than their element with the maximum depth.
% In addition, we allow for the type depth function $\depth$ to be applied directly to terms with the intuitive semantics: $\depth(t) = \depth(\typof(t))$, where $t \in \Terms$.
% Since $\Unknown$ is not constrained, we set $\depth(\Unknown) = \infty$.

\subsection{Precision Order}

The precision order $\preq$ on $\Types$ is defined as the smallest preorder satisfying:
\begin{enumerate}
  \item $\forall\ty \in \Types{:}\{\Unknown\} \leqT \ty$,
  \item $\ty \leqT \ty$ (reflexivity),
  \item $\{\Array{\ty}\} \leqT \{\Array{\ty'}\}$ iff $\ty \leqT \ty'$.
  \item $\{\RegoSet{\ty}\} \leqT \{\RegoSet{\ty'}\}$ iff $\ty \leqT \ty'$.
  \item $(\ty_1 \cup \cdots \cup \ty_m) \leqT (\ty_1' \cup \cdots \cup \ty_n')$ iff $\forall_{1 \leq i \leq n}{:} \exists_{1 \leq j \leq m}{:} \ty_j \preq \ty_i'$, \label{i:uni}
  \item $\{\Obj{v_i{:}\ty_i}\} \leqT \{\Obj{v_i{:}\ty_i'}\}$
     iff for all $i$: $\ty_i \preq \ty_i'$. \label{i:obj}
\end{enumerate}
Constraint \ref{i:uni} says that union $u_1$ is less (or equally) precise than $u_2$ iff every member of $u_2$ is more (or equally) precise than some member of $u_1$.
Constraint \ref{i:obj} describes that for objects $o_1,o_2$ is the relation $o_1 \preq o_2$ true iff both objects have the exact same keys, and for each of their keys $k$, $o_1[k] \preq o_2[k]$.
We additionally write $\ty \ltT \ty'$ for $\ty \leqT \ty'$ and $\ty' \not\leqT \ty$ (strict precision).

The precision order formalises the intuition that \emph{more information is better}: $\ty \leqT \ty'$ reads ``$\ty$ is at most as informative as $\ty'$''.
It is important to note that we have more information about a value which set of possible types is smaller.
The placeholder type $\Unknown$ is the least informative.
A union with fewer alternatives is more informative than one with more.

% \subsection{Union Normalization}
%
% Union types are subject to a \emph{normalization} operation that keeps them in a standard form.
%
% Given a type $\ty$, we define the result of its normalization $\norm(\ty)$ as follows:
% \begin{enumerate}
%   \item If $\ty$ is not a union, return $\ty$ unchanged.
%   \item \emph{Flatten}: replace every union member that is itself a union by the transitive closure of its members.
%   \item \emph{Deduplicate}: remove equal members (by structural equality), keeping first occurrences.
%   \item \emph{Drop unknowns}: if the resulting set has more than one member, remove all $\Unknown$ members.
%   \item \emph{Simplify}: if exactly one member remains, return that member (not a union).
% \end{enumerate}
%
% The effect of $\norm$ is that the empty union is the unique representation for ``no information'', and single-member unions are collapsed.
% After normalization, a union $\ty_1 \cup \cdots \cup \ty_m$ has $m \geq 2$ and no $\Unknown$ member (unless it is the sole member of the whole type, which would have been simplified away).

%%=============================================================
\section{Type Inference Rules}
\label{sec:inference}
%%=============================================================

For the type inference, we use a mutable type map $\typof: \Terms \to \Types$ in accordance with the type function defined in the Section \ref{ss:types}.
We declare a function on $\typof$ giving a set of all already defined terms, noted with $\dom(\typof)$.
An update of the type map at key $v$ with type $\ty$ is noted $\assign{v}{\ty}$; based on it, the two update operations used for the type analysis:

\paragraph{Precision-based update.}
Precision-based update (also called \emph{type refinement}) $\refine{v}{\ty}$ overwrites the existing entry for $v$ with $\ty$ with respect to their precision ordering:
\[
\refine{v}{\ty} =
\begin{cases}
  \typof[v \mapsto \ty] & \text{if } \typof(v) \preq \ty \\
  % \typof[v \mapsto \ty' \cup \ty] & \text{if } \typof(v) = \ty' \cup \ty'' \land \forall t \in \ty'': \{t\} \preq \ty \land \forall t \in \ty': \{t\} \not\preq \ty\\
  \{\Undef\} & \text{otherwise}
\end{cases}
\]

\paragraph{Union extension.}
$\unionize{v}{\ty}$ unionizes the type of $v$ with $\ty$, whilst removing the $\Unknown$ type:
\[
  \unionize{v}{\ty} = \typof\bigl[v \mapsto (\typof(v) \setminus \{\Unknown\}) \cup \ty)\bigr]
\]

\subsection{Term Inference}
\label{sec:term-inference}

\paragraph{Caching.}
Before case analysis, the system checks the cache:
if $v \in \dom(\typof)$ and $\typof(v) \neq \Unknown$, the cached type is returned directly.


\paragraph{Inference of Constants.}
Given the sets $\mathcal{Z}$, containing all whole numbers, and $\Sigma^*$, containing all string literals, the inference for constants of atomic types is straightforward:
\[
\textbf{(T-Null)}\quad
\inferrule{ }{\refine{\Null}{\Null}}
\qquad
\textbf{(T-String)}\quad
\inferrule{ s \in \Sigma^* }{\refine{s}{\Str}}
\]

\[
\textbf{(T-Number)}\quad
\inferrule{ n \in \mathcal{Z} }{\refine{n}{\Int}}
\qquad
\textbf{(T-Bool)}\quad
\inferrule{ b \in \{\top,\bot\} }{\refine{b}{\Bool}}
\]

For the type inference of collections, we have to deliberately take their empty versions into account.

\[
\textbf{(T-ArrayEmpty)}\quad
\inferrule{ }{\refine{[]}{\Array{\Undef}}}
\qquad
\textbf{(T-SetEmpty)}\quad
\inferrule{ }{\refine{\texttt{set()}}{\RegoSet{\Undef}}}
\]

\[
\textbf{(T-Array)}\quad
\inferrule{
  \typof(v_i) = \ty_i \quad \forall i \in 0..n
}{
  \refine{[v_0, \ldots, v_n]}{\Array{\ty_0 \cup \cdots \cup \ty_n}}
}
\]

\[
\qquad
\textbf{(T-Set)}\quad
\inferrule{
  \typof(v_i) = \ty_i \quad \forall i \in 0..n
}{
  \refine{\{v_0, \ldots, v_n\}}{\RegoSet{\ty_0 \cup \cdots \cup \ty_n}}
}
\]

\[
\textbf{(T-Object)}\quad
\inferrule{
  \typof(v_i) = \ty_i \quad \forall i \in 0..n
}{
  \refine{\{k_0{:}v_0, \ldots, k_n{:}v_n\}}{\Obj{k_0{:}\ty_0, \ldots, k_n{:}\ty_n}}
}
\]

% TODO: sig
% if all types match the signature, the type of the whole term is expected
% otherwise we need to unify with undef
\paragraph{Function Call Inference.}
Inferring the types for function calls is described with the rules presented bellow.
\textbf{(T-Call)} applies when there exists a signature of function $f$, that fully matches the provided arguments, i.e., there is no possibility that some argument would be of unexpected type.
In this case, the type of this term is set to the return type $\ty_r$ provided in the signature.

\textbf{(T-CallU)} defines the type of function call to be undefined in a case, when no possible signature could possibly be satisfied by the arguments.
The rule \textbf{(T-CallU)} combines both possibilities when their conditions are not explicitly satisfied.

\[
\textbf{(T-Call)}\quad
\inferrule{
  (\ty_1, \ldots, \ty_k, \ty_r) \in \mathsf{sig}(f) \\
  \typof(t_i) \subseteq \ty_i \quad \forall i \in 1..k
}{
  \refine{f(t_1,\ldots,t_k)}{\ty_r}
}
\qquad
\textbf{(T-CallU)}\quad
\inferrule{
  \typof(t_i) \not\subseteq \ty_i \quad
  \begin{array}{l}
    \forall (\ty_1, \ldots, \ty_k, \ty_r) \in \sig(f) \\
    \exists i \in 1..k
  \end{array}
}{
  \refine{f(t_1,\ldots,t_k)}{\Undef}
}
\]

\[
\textbf{(T-CallDef)}\quad
\inferrule{
  \text{(otherwise)}
}{
  \refine{f(t_1,\ldots,t_k)}{\ty_r \cup \Undef}
}
\]

\paragraph{Reference Inference}

In section \ref{ss:ref}, references of the form \texttt{data.$\pkg$.x} were described to have a value equal to $\sigma(x)$ for global variable $x$ defined in package $\pkg$.
This behaviour is defined with the \textbf{T-PackageRef} rule of inference.
Since the analysis proposed by this paper is restricted only to a fragment of policies contained in a single package, we keep the $\dom$ independent of package names.

\[
\textbf{(T-PackageRef)}\quad
\inferrule{
  x \in \Globals
}{
  \begin{array}{@{}c@{}}
    \refine{\texttt{data.}\pkg.x}{\typof(x)} \\
    \refine{x}{\typof(\texttt{data.}\pkg.x)}
  \end{array}
}
\]

\[
\textbf{(T-ArrM)}\quad
\inferrule{
  x[s_1]\cdots[s_k] \\
  \Array{\ty} \subseteq \typof(x\cdots[s_i]) \\
  i \in 0..(k-1)
}{
  \unionize{x\cdots[s_{i+1}]}{\ty \cup \Undef}
}
\qquad
\textbf{(T-SetM)}\quad
\inferrule{
  x[s_1]\cdots[s_k] \\
  \RegoSet{\ty_s} \subseteq \typof(x\cdots[s_i]) \\
  \typof(s_{i+1}) = \ty \\
  i \in 0..(k-1)
}{
  \unionize{x\cdots[s_{i+1}]}{\ty \cup \Undef}
}
\]

The inference rules \textbf{(T-ArrM)} and \textbf{(T-SetM)} cover accessing members of arrays and sets.
In both cases, we cannot derive solely from the composite type it contains specified key $s_{i+1}$, thus we unionize it with $\Undef$.
For arrays, proposed rule intuitively says that for any \emph{prefix} in reference having an array type, its element corresponds to it.
Since accessing set gives the value of given key (if it is contained in the set), we can derive its value from itself.
The mark $*$ in both rules represents that they are valid even if marked type appears in a union.

\[
\textbf{(T-ObjConM)}\quad
\inferrule{
  r = x[s_1]\cdots[s_k] \\
  \Obj{s_{i+1}{:} \ty, \ldots} \subseteq \typof(x\cdots[s_i]) \\
  i \in 0..(k-1)
}{
  \unionize{x\cdots[s_{i+1}]}{\ty}
}
\]

\vspace{1mm}
%% The case of missing member is covered by the Rego compiler
%%% - no, because this can be used for unions: object{t1} u t2, where a \notin t1 means that accessor of this value on "a" might be undef without crashing the compiler
\[
\textbf{(T-ObjMiss)}\quad
\inferrule{
  r = x[s_1]\cdots[s_k] \\
  \Obj{v_1{:}\ty_1,\ldots,v_n{:}\ty_n} \subseteq \typof(x\cdots[s_i]) \\
  s_{i+1} \notin \bigcup_{i \in 1..n}\{v_i\} \\
  s_{i+1} \notin \Vars
}{
  \unionize{x\cdots[s_{i+1}]}{\Undef}
}
\]

\vspace{1mm}

\[
\textbf{(T-ObjVarM)}\quad
\inferrule{
  r = x[s_1]\cdots[s_k] \\
  \Obj{v_1{:}\ty_1,\ldots,v_n{:}\ty_n} \subseteq \typof(x\cdots[s_i]) \\
  s_{i+1} \in \Vars
}{
  \unionize{x\cdots[s_{i+1}]}{\ty_1 \cup \cdots \cup \ty_n \cup \Undef}
}
\]

% TODO hard to understand
For object types, \textbf{(T-ObjConM)} defines that for an object $o$ and a constant key $k$, accessing $o[k]$ yields corresponding value if $o$ contains given key, whereas \textbf{(T-ObjMiss)} infers such reference $\Undef$ if $o$ does not have a field with key $k$.
When accessing object $o$ with a variable key (inference rule \textbf{(T-ObjVarM)}), we need to take into account every possible field of $o$ and also the case where the key is not present, hence the union.

\subsection{Expression Inference}
\label{sec:expr-inference}

As discussed in \ref{s:terms}, expressions can be either single terms, or function calls, which can themselves be interpreted as terms.
Hence, rules of inference regarding expressions are defined separately to cover the assignment of local variables, which are the cornerstone of type propagation.

\[
\textbf{(E-Assign)}\quad
\inferrule{
  t_L \mathrel{:= } t_R \\
  \typof(t_R) = \ty_R
}{
  \refine{t_L}{\ty_R}
}
\]

\[
\textbf{(E-CallAssign)}\quad
\inferrule{
  f(t_1,\ldots,t_k,t_r) \\
  \typof(f(t_1,\ldots,t_k)) = \ty
}{
  \refine{t_r}{\ty}
}
\]

% Rule \textbf{(E-EqCheck)} handles the pure structural equality test \texttt{==}: both operands are visited for type information, but no type is propagated between them.
% Rule \textbf{(E-Eq)} covers unification/assignment (\texttt{=} and \texttt{:=}) and is the primary mechanism for bidirectional type propagation between variables.
% When a variable $x$ is equated with a term $t$ of type $\ty$, the type $\ty$ is \emph{unioned into} the type of $x$ (and vice versa if both sides are variables).
% This reflects Rego's unification semantics: the equality statement does not simply assign but \emph{constrains}, and both sides carry information about each other.
%
% %% TODO: soft inference
% %% TODO: detection of local/global variables
%% TODO: ranges !!
%% TODO: rules
